\documentclass[11pt]{article}
% Shared local declarations. Document roots use the repository's business-admin family.
\usepackage{business-admin}
\usepackage{amsmath,amssymb,tabularx,longtable,array,booktabs,enumitem}
\usepackage{ragged2e}
\newcolumntype{Y}{>{\RaggedRight\arraybackslash}X}
\newcommand{\SuiteMeta}[2]{\def\DocTitle{#1}\def\DocSubtitle{#2}\def\DocVersion{1.0}\def\DocStatus{Educational and operational reference}\def\DocOwner{Parliamentary Procedure and Governance Suite}\def\DocSummary{Independently authored reference and working materials, based on RONR 12th edition.}}
\newcommand{\SuiteOpen}{\makedoctitle\vspace{-0.5em}\noindent\rule{\linewidth}{0.4pt}\par\medskip}
\newcommand{\RonrAnchor}[1]{\par\smallskip{\footnotesize\textbf{Primary reading:} Robert's Rules of Order Newly Revised, 12th ed., #1. Original synthesis; not official interpretation.}\par}
\newcommand{\WriteLine}{\par\noindent\rule{\linewidth}{0.35pt}\par\vspace{0.5cm}}
\newcommand{\WriteLines}[1]{\par\noindent\rule{\linewidth}{0.35pt}\par\vspace{0.35cm}\noindent\rule{\linewidth}{0.35pt}\par\vspace{0.35cm}\noindent\rule{\linewidth}{0.35pt}\par\vspace{0.25cm}}
\newcommand{\Field}[1]{\par\medskip\noindent\textbf{#1}\par\WriteLines{3}}
\newcommand{\GovCheckbox}{\(\square\)}
\setlength{\emergencystretch}{3em}

\SuiteMeta{Motion Characteristics and Precedence Matrix}{Quick reference | contextual exceptions}
\begin{document}
\SuiteOpen
\RonrAnchor{Sections 5--7, 10--37, and Charts, Tables, and Lists}
\section{Use and limitations}
This quick matrix lists typical characteristics for ordinary circumstances under the 12th edition. It is not a substitute for the Standard Descriptive Characteristics in the cited sections. Special motions, notice, presence, circumstances and adopted rules may change an entry. “Yes” for a second means ordinarily required.

\small\begin{longtable}{@{}p{.22\linewidth}p{.13\linewidth}p{.07\linewidth}p{.14\linewidth}p{.08\linewidth}p{.25\linewidth}@{}}\toprule Motion & Class & 2nd & Debate & Amend & Usual disposition \\ \midrule\endhead
Original main & Main & Yes & Yes & Yes & Majority\\
Postpone indefinitely & Subsidiary & Yes & Yes & No & Majority\\
Amend & Subsidiary & Yes & If pending question debatable & Yes & Majority\\
Commit or refer & Subsidiary & Yes & If pending question debatable & Yes & Majority\\
Postpone definitely & Subsidiary & Yes & Yes (limits) & Yes & Majority; special-order exception\\
Limit/extend debate & Subsidiary & Yes & No & Yes & Two-thirds\\
Previous question & Subsidiary & Yes & No & No & Two-thirds\\
Lay on table & Subsidiary & Yes & No & No & Majority\\
Call for orders & Privileged & No & No & No & Demand; exceptions\\
Question of privilege & Privileged & No & No & No & Chair decides initially\\
Recess & Privileged & Yes & No & Yes & Majority\\
Adjourn & Privileged & Yes & No & No & Majority\\
Fix time to adjourn & Privileged & Yes & No & Yes & Majority\\
Point of order & Incidental & No & No & No & Chair ruling; appeal possible\\
Appeal & Incidental & Yes & Sometimes & No & Majority or tie sustains chair\\
Suspend rules & Incidental & Yes & No & No & Usually two-thirds; limits\\
Object to consideration & Incidental & No & No & No & Two-thirds against consideration\\
Divide question & Incidental & No & No & Yes & Varies with right/decision\\
Division of assembly & Incidental & No & No & No & Demand (not voted)\\
Take from table & Bring back & Yes & No & No & Majority\\
Rescind/amend previous & Bring back & Yes & Yes & Yes & Notice: majority; else two-thirds or entire-membership majority\\
Discharge committee & Bring back & Yes & Yes & Yes & Notice: majority; else two-thirds or entire-membership majority\\
Reconsider & Bring back & Yes & Depends on underlying & No & Majority; mover and timing limits\\
\bottomrule\end{longtable}\normalsize
\section{Source traceability}
Motion names map respectively to RONR (12th ed.) Sections 10--37. For exceptions and reconsideration status consult the original descriptions and the book’s Charts, Tables, and Lists.
\section{Real-time use}
\begin{enumerate}[leftmargin=*,itemsep=4pt]
\item Identify the pending question and the class of the proposed motion.
\item Check whether motion is in order and whether the member may interrupt.
\item Inspect the full section for contextual qualifications.
\item Confirm second, debate, amendment, notice and denominator; announce the result.
\end{enumerate}
\end{document}
